% !TEX program = xelatex
\documentclass[12pt,a4paper]{report}

% 1. Pengaturan Font & Bahasa (Wajib XeLaTeX / LuaLaTeX)
\usepackage{fontspec}
\usepackage{polyglossia}
\usepackage{libertine}

% Bahasa utama dokumen (Indonesia)
\setdefaultlanguage{indonesian}
% Bahasa kedua (Arab)
\setotherlanguage{arabic}

\newfontfamily\arabicfont[Script=Arabic, Scale=0.8]{Amiri}

% 2. Pengaturan Halaman & Spasi
\usepackage[margin=2.5cm]{geometry}
\usepackage{setspace}
\onehalfspacing

% 3. Bingkai Matan ala Kitab Syarah
\usepackage{tcolorbox}
\usepackage{xcolor}

\definecolor{matanbg}{RGB}{250, 248, 242}
\definecolor{matanborder}{RGB}{120, 90, 40}

\newtcolorbox{matanbox}[1][Matan]{
  colback=matanbg,
  colframe=matanborder,
  boxrule=0.8pt,
  arc=2pt,
  left=8pt,
  right=8pt,
  top=6pt,
  bottom=6pt,
  fonttitle=\bfseries\small,
  title={#1},
  before upper={\parindent=0pt }
}

% Penyesuaian penamaan BAB
\addto\captionsindonesian{
  \renewcommand{\chaptername}{BAB}
}

\hyphenpenalty=500
\exhyphenpenalty=500
\emergencystretch=2em

\begin{document}

% Membuat halaman judul atau daftar isi
\tableofcontents
\newpage

% Memanggil file bab terpisah (tanpa ekstensi .tex)
\chapter*{Pengantar}

\begin{center}
\textarabic{﷽}
\end{center}

Kepada-Nya aku bersandar, kepada-Nya rasa \textit{tsiqah} percayaku berlabuh, dan kepadanya pula aku memohon pertolongan atas seluruh urusan dunia dan agama.

Segala puji bagi Allah yang telah menunjukkan kita untuk mengenal yang halal dan haram, dan mewajibkan atas kita untuk taat kepada Nabi kita Muhammad, pemimpin seluruh makhluk. Dia juga memerintahkan kita untuk mengikuti syariat-Nya yang agung dan memahami hukum-hukum, serta memperbolehkan kita memandang Wajah-Nya Yang Maha Mulia di Darussalam. Dia Allah membenci bagi kita kekufuran, kefasikan, dan kemaksiatan, serta mengharamkan atas kita perbuatan dosa dan kesewenang-wenangan.

Aku bersaksi bahwa tidak ada \textit{illah} yang berhak diibadahi selain Allah semata, tiada sekutu bagi-Nya. Pemilik keagungan, kemuliaan, keperkasaan, dan kekuasaan. Aku bersaksi bahwa Muhammad adalah hamba dan utusan-Nya, sang penyeru kepada pendalaman fikih agama. Semoga shalawat dan salam Allah tercurahkan kepadanya, keluarganya, dan para sahabatnya sekalian, \textit{amma ba’du}.

Risalah ini adalah buku ringkas dalam memilih pendapat-pendapat paling kuat dari para ulama berkaitan dengan sebagian permasalahan-permasalahan thaharah dari kitab matan \textbf{Safinatun Najah} dan \textbf{Nailur Raja’ Bisyarhi Safinatun Naja}. Buku ini dikhususkan bagi santriwati setingkat SMP. Saya menamai buku ini dengan judul \textbf{Syarhu Safinatin Najati Fi Ba’dhi Ahkami Thaharati Li Thalibat Mustawa Tsanawi}.

Saya menyusunnya sebagai pendekatan pengantar bagi pemula dengan mencukupkan pada isyarat terhadap dalil, tanpa menyelami rincian, bantahan serta uraian yang panjang. Masalah-masalah fiqih yang dibahas merupakan pendapat-pendapat dari madzhab Syafi’i \textit{rahimahullah}, yang diambil dari pendapat beliau, pilihan ulama pengikut beliau dan ulama madzhab yang masyhur. Dalam mengambil pendapat, saya tidak terikat dengan pendapat masyhur dari ulama \textit{muta’akhirin} (kontemporer), akan tetapi pegangan dalam hal ini mengacu pada pendapat yang memiliki dalil, berdasarkan penelitian Syeikh Abu Qatadah Al-filistini \textit{hafidzahullah} dalam kitabnya \textbf{Mukhtashar Abi Qatadah Fi Fiqhi Asy-Syafi'iyah As-Sadah}.

Syeikh Syeikh Abu Qatadah Al-filistini \textit{hafidzahullah} mengatakan: 

"Berkomitmen pada suatu madzhab fikih bukan berarti menolak dalil dan tarjih (memilih pendapat yang paling kuat). Madzhab-madzhab ulama hakikatnya diletakkan untuk mendekatkan terhadap dalil dan penjelasan terhadap pendapat imam madzhab, para pengikutnya, dan \textit{ashabul wujuh } (ulama yang mengembangkan madzhab) bukan supaya taklid kecuali bagi orang awam yang keadaannya tidak memungkinkan untuk selain itu.Kemudian tidak ada suatu madzhab kecuali di dalamnya terdapat ruang yang luas untuk berijtihad dan memilih pendapat. Inilah tujuan saya setelah mempermudah masalah-masalah fikih dan meringkasnya, dan Allahlah yang memberikan taufik.

Bagi siapa yang membaca dan menelaah madzhab Syafi’i akan mengetahui keluasannya, dan sesungguhnya para ulama madzhab tidaklah mengambil pendapat Imam Syafi’i kecuali memiliki dasar dalil untuk berijtihad dan memilih pendapat bukan untuk taqlid yang justru dibenci oleh imam Syafi’i itu sendiri yang telah melarang hal tersebut seperti yang disebutkan oleh Immam  Al-Muzani \textit{rahimahullah} dalam Mukhtasarnya. Maka barang siapa yang bertaqlid kepada Syafi’i tanpa penelitian, ilmu, dan pengetahuan akan dalilnya sesungguhnya dia telah membatalkan tujuan utama Imam Syafi’i dalam ilmu ini. Hal ini sepantasnya harus dipahami dalam perkara ini. Tidak diperbolehkan bagi seseorang untuk menolak pendapat orang lain dengan alasan bahwa pendapat yang ia pilih tersebut menyelisihi pendapat yang masyhur dalam madzhab. Padahal pendapat yang masyhur itu sebenarnya hanyalah ijtihad satu atau lebih para ulama. Sedangkan dia hanya sedang memilih suatu pendapat dari para ulama -Selesai perkataan.

Allah tempat meminta, agar menjadikan karya ini ikhlas semata-mata hanya karena mengharapkan wajah Allah yang mulia, dan menjadikan kitab ini tersebar luas seperti peredaran matahari di penjuru ufuk, serta memberikan manfaat pada penulisnya dan para penghafalnya, dan orang yang menyibukkan diri dengannya. Tiada taufik bagiku kecuali dengan pertolongan dari Allah, hanya kepada-Nya aku bertawakal dan bertaubat. Dialah Allah sebaik-baik pelindung dan penolong. Maha suci Rabbmu, Rabb pemilik kemuliaan dari apa yang mereka sifatkan, dan salam atas seluruh rasul, \textit{walhamdulillahi rabbil 'alamin}.

\vspace{1cm} % Memberi jarak vertikal antara teks dan blok penulis

\begin{flushleft}
Ditulis oleh, hamba yang paling fakir pada Allah \textit{ta'ala}\\
\textbf{Zen Ibrahim} semoga Allah mengampuninya dan kedua orang tuanya.\\
\textbf{Markaz Mulazamah Mutiara Muslimah}\\
Karanganyar - Jawa Tengah\\
Hari Kamis 14 Syawal 1447 H bertepatan dengan 2 April 2026 M
\end{flushleft}
\chapter{Kitab Thaharah}

\section{Definisi Thaharah}

\textbf{Makna kitab thaharah:} Kumpulan hukum-hukum thaharah.

\textbf{Thaharah secara syar'i:} Mengangkat hadats atau menghilangkan najis, dan apa-apa yang semakna dengan keduanya atau yang serupa dengan keduanya.

\textbf{Najis:} Suatu kotoran yang menjijikkan menghalangi sahnya shalat, selama tidak ada rukhshah; yaitu keringanan yang membolehkannya, seperti kotoran tinja manusia itu sendiri.

\textbf{Mutananjis:} Sesuatu yang suci yang terkena najis seperti pakaian dan karpet.

\textbf{Mengangkat hadats:} Menghilangkan sifat syar'i yang melekat pada badan yang menghalangi sahnya shalat atau semisalnya. Contohnya: Wudhu untuk mengangkat hadats kecil dan mandi untuk mengangkat hadats besar. 

Hadats dibagi menjadi 2 jenis:
\begin{enumerate} 
    \item \textit{Hadats asghar} (hadats kecil): yaitu hadats yang diangkat wudhu.
    \item \textit{Hadats akbar} (hadats besar): yaitu hadats yang diangkat dengan mandi.
\end{enumerate}

Maksud dari \textit{"dan apa-apa yang semakna dengan keduanya"}; yaitu apa-apa yang semakna dengan mengangkat hadats dan yang semakna dengan menghilangkan najis.

\begin{enumerate}
    \item Apa-apa yang semakna dengan mengangkat hadats:
    \begin{itemize}
        \item Memandikan jenazah. Karena memandikannya tidak disebabkan oleh hadats, tetapi sebagai perkara \textit{taabud} (sebagai bentuk ibadah) dan perbuatan itu tidak mengangkat hadats.
        \item Tayamum. Dalam madzhab, tayamum tidak mengangkat hadats, melainkan hanya membolehkan shalat dan yang semisalnya yang diharamkan sebab hadats. Oleh karena itu mereka mengatakan: Tayamum bersifat \textit{mubih} (membolehkan) dan bukan mengangkat hadats. Sedangkan dalam Mukhtasar Abi Qatadah Fi Fiqhi Asy-Syafi’i As-Sadah: Tayamum seperti wudhu yang mengangkat hadats.
    \end{itemize} 
    \item Apa-apa yang semakna dengan menghilangkan najis: Seperti mengusap, menggosok dan dengan sinar matahari.
\end{enumerate}
	
Maksud dari \textit{yang serupa dengan keduanya}; yaitu serupa dalam mengangkat hadats dan serupa dalam menghilangkan najis:
\begin{enumerate}
    \item Yang serupa dalam mengangkat hadats:
    \begin{itemize}
        \item Basuhan kedua, dan ketiga dalam berwudhu. Hadats sebenarnya telah terangkat pada basuhan pertama. Adapun basuhan kedua dan ketiga adalah sunah inilah yang disebut \textit{yang serupa dengan mengangkat hadats} yang pada hakikatnya tidak mengangkat hadats.
        \item Memperbarui wudhu.
        \item Mandi yang disunahkan seperti mandi jum’at.
    \end{itemize}
    \item Yang serupa dalam menghilangkan najis: Seperti basuhan kedua dan ketiga pada tempat yang mengeluarkan najis (dubur dan qubul) setelah kita menghilangkan zat najisnya.
\end{enumerate}

\textbf{Hukum Thaharah:} Bersuci dari najis dan menghilangkannya hukumnya wajib bagi yang ingat dan mampu. Adapun bersuci dari hadats hukumnya wajib untuk membolehkan shalat.

\section{Tujuan, Sarana, dan Perantara Sarana Thaharah}

Tujuan thaharah ada empat: 
\begin{enumerate}
    \item Wudhu.
    \item Mandi. 
    \item Tayamum. 
    \item Menghilangkan najis.
\end{enumerate}

Sarana thaharah ada empat: 
\begin{enumerate}
    \item Air.
    \item Debu.
    \item Batu. 
    \item Menyamak kulit.
\end{enumerate}

Perantara sarana thaharah ada dua: 
\begin{enumerate}
    \item Bejana.
    \item Ijitihad.
\end{enumerate}

\chapter{Tanda-tanda Baligh}

\begin{matanbox}[Matan]
Tanda-tanda baligh ada tiga:
\begin{enumerate}
    \item Mencapai genap umur 15 tahun bagi laki–laki dan perempuan.
    \item \textit{Ihtilam} (mimpi basah) pada umur 9 tahun bagi laki–laki dan perempuan.
    \item Haid pada usia 9 tahun bagi perempuan.
\end{enumerate}
\end{matanbox}

\textbf{Tanda-tanda (\textarabic{اَلْعَلَامَاتُ}):} Merupakan bentuk jamak dari kata \textit{'alamah} (\textarabic{عَلَامَة}) artinya petunjuk.

\textbf{Tanda (\textarabic{عَلَامَة}) menurut istilah syariat adalah:} Suatu yang dari keberadaannya berkonsekuensi adanya hukum, namun dari ketiadaannya tidak berkonsekuensi ketiadaan hukum.

\textbf{Maksudnya:} Adanya \textit{ihtilam} berkonsekuensi adanya baligh, namun bukan berarti jika tidak adanya \textit{ihtilam} mesti berkonsekuensi tidak adanya baligh. Karena bisa jadi dia baligh bukan karena tanda \textit{ihtilam}; bisa jadi baligh karena tanda haid, dan bisa jadi karena tanda genap berusia lima belas tahun.

\textbf{Bulugh atau baligh}: Telah sampainya pada batasan \textit{taklif} (pembebanan syariat), dan disyaratkan berakal bersamanya.

\section{Mencapai genap umur 15 tahun bagi laki–laki dan perempuan}

\textbf{Maksudnya:} Berlalunya 15 tahun perhitungan kalendar \textit{qamariyah tahdidiyah} atas \textit{shabiyi} atau \textit{shabiyah} sejak kelahiran sempurnanya.

\textbf{Makna shabiyi:} Usia anak laki-laki dari lahir sampai baligh.

\textbf{Makna shabiyah:} Usia anak perempuan dari lahir sampai baligh.

\textbf{Makna qamariyah tahdidiyah}: Penetapan usia secara tepat dengan menggunakan perhitungan kalender qamariyah yang ditentukan secara tepat.

\section{Ihtilam pada umur 9 tahun bagi laki–laki dan perempuan}

\textbf{Ihtilam:} Apa yang dimimpikan oleh seseorang dalam tidur.
\textbf{Maksudnya:} Keluarnya air mani setelah sempurna berusia 9 tahun dengan perhitungan kalender \textit{qamariyah taqribiyah}.
\textbf{Makna Qamariyah Taqribiyah:} Penetapan usia secara perkiraan dengan menggunakan perhitungan kalender qamariyah tanpa terikat hitungan hari secara persis.

\section{Haid pada usia 9 tahun bagi perempuan}

\textbf{Haid secara istilah syariat:} Darah alami yang keluar dari bagian pangkal terdalam rahim perempuan, dalam keadaan sehat, pada waktu tertentu.

\textbf{Maksudnya:} Perempuan mengalami haid setelah mencapai usia 9 tahun melalui perhitungan kalender \textit{qamariah tahdidiyah}, tidaklah mengapa jika saat haid usianya belum tepat 9 tahun, sekurangnya adalah waktu yang tidak cukup untuk masa haid dan suci; yaitu tidak kurang dari 16 hari. 

% Contoh penggunaan di dalam dokumen utama atau file bab (misal: kitab_thaharah.tex)


\end{document}